\section{Esperimenti sui livelli di reasoning effort}
\label{sec:effort}

Il commit \texttt{1b6a7e2} (30 giugno 2026, \emph{``ADDED: new notebook to track differences in low,medium,high effort performances''}) introduce \texttt{2\_2\_benchmark\_different\_efforts.ipynb}, che ripete l'esperimento di \S\ref{sec:openai} facendo variare esplicitamente il parametro \texttt{reasoning.effort} di o3-mini su tre livelli (\texttt{low}, \texttt{medium}, \texttt{high}), sia sul dataset Adult sia su un secondo dataset, Student Math, non utilizzato negli esperimenti precedenti. I sei file JSON prodotti (tre per dataset) usano lo stesso prompt di base (3662 token di input, costanti sulle tre run Adult), isolando quindi l'effetto del solo parametro \texttt{effort}.

I risultati su Adult non mostrano un miglioramento monotono dell'accuratezza all'aumentare dello sforzo di ragionamento:

\begin{center}
\begin{tabular}{lrrr}
\toprule
Effetto & \texttt{low} & \texttt{medium} & \texttt{high} \\
\midrule
TV (errore \%) & 6.16 & 6.16 & 6.10 \\
TE (errore \%) & 1.39 & 8.66 & 38.73 \\
SE (errore \%) & 57.51 & 309.70 & 879.98 \\
DE (errore \%) & 35.50 & 12.44 & 9.45 \\
IE (errore \%) & 100.00 & 202.58 & 325.54 \\
\midrule
Tempo (s) & 192.3 & 175.2 & 203.1 \\
Token totali & 68\,753 & 77\,849 & 85\,021 \\
\bottomrule
\end{tabular}
\end{center}

Il pattern è controintuitivo: TE, SE e IE \emph{peggiorano} nettamente passando da \texttt{low} a \texttt{high} (l'errore su SE arriva a quasi il 900\%), mentre solo DE migliora. Il tempo di risposta resta nell'ordine di 3--3.5 minuti indipendentemente dal livello impostato, e il consumo di token cresce senza un corrispondente guadagno di accuratezza complessivo. Il repository non contiene un'analisi esplicita scritta all'epoca di questo commit che ne spieghi la causa; l'osservazione più diretta che se ne può trarre è che aumentare lo sforzo di ragionamento non è, di per sé, una leva efficace per migliorare l'affidabilità del calcolo. Questo risultato, insieme al costo già elevato osservato in \S\ref{sec:analisi-json} anche al livello \texttt{low}, prepara il terreno per la valutazione dei limiti del modello descritta nella fase successiva.
